% ============================================================================
% PREÁMBULO DEL LIBRO
% Curso práctico de Business Intelligence con R y RStudio
% ----------------------------------------------------------------------------
% Compila con pdfLaTeX (Overleaf lo usa por defecto) o con:
%     latexmk -pdf curso-bi-rstudio.tex
% ============================================================================

% ---------------------------------------------------------------------------
% Idioma, codificación y fuentes
% ---------------------------------------------------------------------------
\usepackage[utf8]{inputenc}
\usepackage[T1]{fontenc}
% es-noshorthands: evita que babel active los caracteres " < > (rompen el código)
% es-nolists: conserva las viñetas estándar en las listas
% es-tabla: las tablas se llaman "Tabla" y no "Cuadro"
\usepackage[spanish,es-noshorthands,es-nolists,es-tabla,es-ucroman]{babel}
\usepackage{lmodern}
\usepackage{microtype}
\usepackage{textcomp}

% ---------------------------------------------------------------------------
% Página
% ---------------------------------------------------------------------------
\usepackage[letterpaper,top=2.5cm,bottom=2.5cm,left=2.5cm,right=2.5cm,
            headheight=15pt]{geometry}
\setlength{\parindent}{0pt}
\setlength{\parskip}{0.55em}
\linespread{1.08}

% ---------------------------------------------------------------------------
% Paquetes generales
% ---------------------------------------------------------------------------
\usepackage{xcolor}
\usepackage{graphicx}
\graphicspath{{figuras/}}
\usepackage{amsmath,amssymb}
\usepackage{booktabs,tabularx,array,longtable,multirow}
\usepackage{enumitem}
\setlist{itemsep=0.15em, topsep=0.3em}
\usepackage{etoolbox}
\usepackage{float}
\usepackage[font=small,labelfont={bf,sf},labelsep=period]{caption}
\usepackage{tikz}
\usetikzlibrary{positioning,arrows.meta,shapes.geometric,calc,fit,backgrounds}
\usepackage{listings}
\usepackage[most]{tcolorbox}
\usepackage{fancyhdr}
\usepackage{titlesec}
\usepackage{pifont}
\usepackage{xurl}

% ---------------------------------------------------------------------------
% Paleta de colores del libro
% ---------------------------------------------------------------------------
\definecolor{primario}{RGB}{23,63,105}      % azul marino
\definecolor{secundario}{RGB}{0,121,140}    % azul verdoso
\definecolor{acento}{RGB}{214,110,24}       % naranja
\definecolor{exito}{RGB}{46,125,50}         % verde
\definecolor{peligro}{RGB}{183,28,28}       % rojo
\definecolor{morado}{RGB}{94,53,177}        % morado
\definecolor{gristexto}{RGB}{90,90,90}

% Colores del código
\definecolor{codigofondo}{RGB}{246,248,250}
\definecolor{codigoborde}{RGB}{208,215,222}
\definecolor{salidafondo}{RGB}{255,255,255}
\definecolor{rpalabra}{RGB}{0,92,197}
\definecolor{rcomentario}{RGB}{106,115,125}
\definecolor{rtexto}{RGB}{3,47,98}
\definecolor{rtextocadena}{RGB}{163,21,21}
\definecolor{codigoinline}{RGB}{176,48,96}

% ---------------------------------------------------------------------------
% Títulos de capítulos y secciones
% ---------------------------------------------------------------------------
\addto\captionsspanish{%
  \renewcommand{\chaptername}{Módulo}%
  \renewcommand{\contentsname}{Contenido}%
  \renewcommand{\appendixname}{Apéndice}%
}

\titleformat{\chapter}[display]
  {\normalfont\sffamily\color{primario}}
  {\large\bfseries\MakeUppercase{\chaptertitlename}\ \thechapter}
  {6pt}
  {\Huge\bfseries}
  [{\color{acento}\titlerule[2pt]}]
\titleformat{name=\chapter,numberless}[display]
  {\normalfont\sffamily\color{primario}}
  {}
  {0pt}
  {\Huge\bfseries}
  [{\color{acento}\titlerule[2pt]}]
\titlespacing*{\chapter}{0pt}{-10pt}{24pt}

\titleformat{\part}[display]
  {\normalfont\sffamily\centering\color{primario}}
  {\Large\bfseries\MakeUppercase{\partname}\ \thepart}
  {14pt}
  {\fontsize{30}{36}\selectfont\bfseries}
  [{\vspace{10pt}\color{acento}\titlerule[2pt]}]

\titleformat{\section}
  {\normalfont\sffamily\Large\bfseries\color{primario}}
  {\thesection}{0.8em}{}
\titleformat{\subsection}
  {\normalfont\sffamily\large\bfseries\color{secundario}}
  {\thesubsection}{0.8em}{}
\titleformat{\subsubsection}
  {\normalfont\sffamily\normalsize\bfseries\color{gristexto}}
  {}{0em}{}

\setcounter{secnumdepth}{2}
\setcounter{tocdepth}{1}

% ---------------------------------------------------------------------------
% Encabezados y pies de página
% ---------------------------------------------------------------------------
\pagestyle{fancy}
\fancyhf{}
\makeatletter
\renewcommand{\chaptermark}[1]{%
  \markboth{\if@mainmatter\chaptername\ \thechapter.\ \fi#1}{}}
\makeatother
\fancyhead[L]{\small\sffamily\color{gristexto}\leftmark}
\fancyhead[R]{\small\sffamily\color{gristexto}Business Intelligence con R}
\fancyfoot[C]{\small\sffamily\thepage}
\renewcommand{\headrulewidth}{0.4pt}
\fancypagestyle{plain}{%
  \fancyhf{}%
  \fancyfoot[C]{\small\sffamily\thepage}%
  \renewcommand{\headrulewidth}{0pt}%
}

% ---------------------------------------------------------------------------
% Enlaces
% ---------------------------------------------------------------------------
\usepackage{hyperref}
\hypersetup{
  colorlinks=true,
  linkcolor=primario,
  urlcolor=secundario,
  citecolor=secundario,
  pdftitle={Curso práctico de Business Intelligence con R y RStudio},
  pdfsubject={Business Intelligence, Business Analytics, R, RStudio},
  pdfkeywords={R, RStudio, BI, dplyr, ggplot2, Shiny, machine learning},
  bookmarksnumbered=true
}
\usepackage{bookmark}

% ---------------------------------------------------------------------------
% Estilos de código (listings)
% ---------------------------------------------------------------------------
% Mapeo de caracteres UTF-8 para que los acentos se vean dentro del código
% cuando se compila con pdfLaTeX.
\lstset{
  inputencoding=utf8,
  extendedchars=true,
  literate=
    {á}{{\'a}}1 {é}{{\'e}}1 {í}{{\'i}}1 {ó}{{\'o}}1 {ú}{{\'u}}1
    {Á}{{\'A}}1 {É}{{\'E}}1 {Í}{{\'I}}1 {Ó}{{\'O}}1 {Ú}{{\'U}}1
    {ñ}{{\~n}}1 {Ñ}{{\~N}}1 {ü}{{\"u}}1 {Ü}{{\"U}}1
    {¿}{{\textquestiondown}}1 {¡}{{\textexclamdown}}1
    {°}{{\textdegree}}1 {×}{{$\times$}}1 {…}{{\ldots}}1
    {«}{{\guillemotleft}}1 {»}{{\guillemotright}}1
    {→}{{$\rightarrow$}}1 {✓}{{\ding{51}}}1 {✗}{{\ding{55}}}1
    {²}{{\textsuperscript{2}}}1 {€}{{EUR}}3
}

\lstdefinestyle{estiloR}{
  language=R,
  basicstyle=\ttfamily\small\color{rtexto},
  keywordstyle=\color{rpalabra},
  commentstyle=\color{rcomentario}\itshape,
  stringstyle=\color{rtextocadena},
  showstringspaces=false,
  upquote=true,
  columns=fixed, basewidth=0.52em,
  keepspaces=true,
  breaklines=true,
  breakatwhitespace=false,
  postbreak=\mbox{\textcolor{gristexto}{$\hookrightarrow$}\space},
  tabsize=2,
  numbers=left,
  numberstyle=\tiny\color{gray},
  numbersep=8pt,
  xleftmargin=12pt,
  morekeywords={TRUE,FALSE,NA,NULL,Inf,NaN,function,library,require,
                if,else,for,while,repeat,break,next,return,in},
  deletekeywords={c,t,q,data,labels,model,colors,names,levels,table,
                  text,title,frame,length,list,matrix,format,file,
                  scale,sort,order,rank,rev,trunc,round,range},
}

\lstdefinestyle{estiloSalida}{
  basicstyle=\ttfamily\small\color{black!80},
  columns=fixed, basewidth=0.52em,
  keepspaces=true,
  breaklines=true,
  postbreak=\mbox{\textcolor{gristexto}{$\hookrightarrow$}\space},
  upquote=true,
  numbers=none,
  xleftmargin=0pt,
}

\lstdefinestyle{estiloSQL}{
  language=SQL,
  basicstyle=\ttfamily\small\color{rtexto},
  keywordstyle=\color{rpalabra},
  commentstyle=\color{rcomentario}\itshape,
  stringstyle=\color{rtextocadena},
  showstringspaces=false,
  upquote=true,
  columns=fixed, basewidth=0.52em,
  keepspaces=true,
  breaklines=true,
  morekeywords={LIMIT,OVER,PARTITION,WITH,CASE,WHEN,THEN,END,COALESCE,
                ROUND,SUM,AVG,COUNT,MIN,MAX,DISTINCT,INNER,LEFT,RIGHT,JOIN},
}

\lstdefinestyle{estiloTexto}{
  basicstyle=\ttfamily\small\color{rtexto},
  commentstyle=\color{rcomentario}\itshape,
  columns=fixed, basewidth=0.52em,
  keepspaces=true,
  breaklines=true,
  upquote=true,
  numbers=left,
  numberstyle=\tiny\color{gray},
  numbersep=8pt,
  xleftmargin=12pt,
}

% ---------------------------------------------------------------------------
% Entornos de código
% ---------------------------------------------------------------------------
% rcode       -> código R que se EJECUTA (el libro lo verifica en orden)
% rnoejecutar -> código R solo para mostrar (instalaciones, apps Shiny,
%                conexiones a servidores externos, plantillas)
% rsalida     -> lo que R imprime en la consola
% sqlcode     -> consultas SQL
% textocodigo -> R Markdown, YAML, terminal u otros textos literales
% Todos aceptan opciones de tcolorbox en la misma línea, por ejemplo:
%   \begin{rcode}[title=Mi primer script]
\tcbset{
  basecodigo/.style={
    enhanced, breakable,
    colback=codigofondo, colframe=codigoborde,
    boxrule=0.5pt, arc=2pt,
    left=4pt, right=4pt, top=2pt, bottom=2pt,
    before skip=8pt, after skip=8pt,
    colbacktitle=primario!85, coltitle=white, fonttitle=\sffamily\bfseries\small,
  }
}

\DeclareTCBListing{rcode}{ !O{} }{
  basecodigo, listing only,
  listing options={style=estiloR},
  #1
}

\DeclareTCBListing{rnoejecutar}{ !O{} }{
  basecodigo, listing only,
  listing options={style=estiloR},
  borderline west={2pt}{0pt}{acento!70},
  #1
}

\DeclareTCBListing{rsalida}{ !O{} }{
  enhanced, breakable, listing only,
  colback=salidafondo, colframe=codigoborde,
  boxrule=0.4pt, arc=2pt, leftrule=3pt,
  colframe=secundario!45,
  left=6pt, right=4pt, top=1pt, bottom=1pt,
  before skip=4pt, after skip=10pt,
  listing options={style=estiloSalida},
  #1
}

\DeclareTCBListing{sqlcode}{ !O{} }{
  basecodigo, listing only,
  listing options={style=estiloSQL},
  borderline west={2pt}{0pt}{morado!60},
  #1
}

\DeclareTCBListing{textocodigo}{ !O{} }{
  basecodigo, listing only,
  listing options={style=estiloTexto},
  #1
}

% ---------------------------------------------------------------------------
% Cajas didácticas
% ---------------------------------------------------------------------------
\tcbset{
  basecaja/.style={
    enhanced, breakable,
    boxrule=0.6pt, arc=3pt,
    left=8pt, right=8pt, top=6pt, bottom=6pt,
    before skip=10pt, after skip=10pt,
    fonttitle=\sffamily\bfseries,
    coltitle=white,
  }
}

% Objetivos al inicio de cada módulo
\newtcolorbox{objetivos}[1][]{
  basecaja, colback=primario!4, colframe=primario,
  title={\ding{43}\ \ Al terminar este módulo podrás\ldots}, #1
}

% Concepto clave (teoría importante)
\newtcolorbox{concepto}[2][]{
  basecaja, colback=secundario!5, colframe=secundario,
  title={Concepto clave: #2}, #1
}

% Ejemplo práctico resuelto y explicado
\newtcolorbox{ejemplo}[2][]{
  basecaja, colback=exito!4, colframe=exito,
  title={Ejemplo práctico: #2}, #1
}

% Caso de negocio (aplicación real en una empresa)
\newtcolorbox{negocio}[2][]{
  basecaja, colback=morado!4, colframe=morado,
  title={En la empresa: #2}, #1
}

% Consejo profesional
\newtcolorbox{consejo}[1][]{
  basecaja, colback=secundario!3, colframe=secundario!70,
  title={\ding{46}\ \ Consejo}, #1
}

% Advertencia / error común
\newtcolorbox{advertencia}[1][]{
  basecaja, colback=peligro!4, colframe=peligro!85,
  title={\ding{115}\ \ Cuidado}, #1
}

% Ejercicio para el lector (numerado por módulo).
% Uso: \begin{ejercicio}{m2-3}{Título}  ...  \end{ejercicio}
% El primer argumento es un identificador único; la solución correspondiente
% se escribe en el apéndice con \begin{solucion}{m2-3} ... \end{solucion}.
\newtcolorbox[auto counter, number within=chapter]{ejercicio}[3][]{
  basecaja, colback=acento!5, colframe=acento,
  title={Ejercicio~\thetcbcounter: #3}, label={ej:#2},
  after upper={\par\smallskip\hfill{\small\sffamily\color{gristexto}%
    \ding{43}\ Solución en la página~\pageref{sol:#2}}},
  #1
}

% "Hazlo tú": práctica guiada dentro del texto
\newtcolorbox{hazlotu}[1][]{
  basecaja, colback=acento!3, colframe=acento!70,
  title={\ding{45}\ \ Hazlo tú}, #1
}

% Solución de un ejercicio (se usa en el apéndice de soluciones)
% Uso: \begin{solucion}{m2-3} ... \end{solucion}
\newtcolorbox{solucion}[2][]{
  basecaja, colback=exito!3, colframe=exito!70,
  title={Solución del ejercicio~\ref*{ej:#2}},
  before upper={\phantomsection\label{sol:#2}},
  after upper={\par\smallskip\hfill{\small\sffamily\color{gristexto}%
    Enunciado en la página~\pageref{ej:#2}}},
  #1
}

% Resumen del módulo
\newtcolorbox{resumen}[1][]{
  basecaja, colback=primario!3, colframe=primario!80,
  title={Resumen del módulo}, #1
}

% ---------------------------------------------------------------------------
% Comandos útiles
% ---------------------------------------------------------------------------
% Código en línea: \codigo{mutate()}. Escapa _ % $ # & { } con barra invertida.
\newcommand{\codigo}[1]{\textcolor{codigoinline}{\texttt{#1}}}
% Nombre de paquete: \paquete{dplyr}
\newcommand{\paquete}[1]{\textsf{\textbf{#1}}}
% Nombre de archivo o carpeta: \archivo{ventas.csv}
\newcommand{\archivo}[1]{\texttt{#1}}
% Opción de menú de RStudio: \menu{Tools > Global Options}
\newcommand{\menu}[1]{\textsf{\textbf{\color{primario}#1}}}
% Tecla o atajo de teclado: \tecla{Ctrl} + \tecla{Enter}
\newtcbox{\tecla}{on line, boxsep=0pt, left=3pt, right=3pt, top=1.5pt,
  bottom=1.5pt, colback=gray!8, colframe=gray!55, boxrule=0.5pt, arc=2pt,
  fontupper=\ttfamily\footnotesize}
% Nivel de dificultad de un ejercicio: \nivel{Básico}
\newcommand{\nivel}[1]{{\small\sffamily\color{gristexto}Nivel: \textbf{#1}}\par}
% Término nuevo resaltado la primera vez que aparece
\newcommand{\termino}[1]{\textbf{\color{primario}#1}}
